\documentclass[11pt]{article}
\usepackage[margin=1in]{geometry}
\usepackage{amsmath,amssymb,amsthm}
\usepackage[T1]{fontenc}
\usepackage{lmodern}
\usepackage{xcolor}
\usepackage{listings}
\usepackage{booktabs}
\usepackage{longtable}
\usepackage{array}
\usepackage[hidelinks]{hyperref}
\usepackage{url}

\newtheorem{theorem}{Theorem}
\newtheorem{lemma}[theorem]{Lemma}
\newtheorem{corollary}[theorem]{Corollary}
\newtheorem{proposition}[theorem]{Proposition}
\theoremstyle{definition}
\newtheorem{definition}[theorem]{Definition}
\theoremstyle{remark}
\newtheorem{remark}[theorem]{Remark}

% \lean{...}: a Lean identifier. Every argument of \lean is checked by the name gate
% (gate_names.sh) against the five source trees and the pinned Mathlib.
\makeatletter
\def\UrlBreaks{\do\.\do\_\do\'}
\makeatother
\DeclareUrlCommand\leanu{\def\UrlFont{\ttfamily\upshape}}
\DeclareRobustCommand\lean[1]{\ifmmode\text{\leanu{#1}}\else\leanu{#1}\fi}
% \code{...}: Lean text that is not a single identifier (not name-gated).
\newcommand{\code}[1]{\texttt{#1}}
% \ext{...}: a Lean name from a third-party development (breakable, not name-gated).
\DeclareUrlCommand\ext{\def\UrlFont{\ttfamily\upshape}}

\newcommand{\RR}{\mathbb{R}}

% version-2 commits: carnot=fa8c37951d290d26db66ff83c358b09b156d26e6 viviani=e71460f2edb224b2e7e98a39c6db6bead93613cc degua=daa89522f4a85e6afb8fe4d4d1a94b3d0761943e japanese-theorem=be4503b4e1dfdf7e527e2425276f992ecabb4bc2 poncelet-triangles=aa9da6c4e077cee2e9191ad9011a246ad9e04f2f
\newcommand{\commitcarnot}{98e5614812f2fc2491a729adf9677bfce9b55217}
\newcommand{\commitviviani}{029aa3d464db2ac8b3eb2e8f57e79e80f5585803}
\newcommand{\commitdegua}{d46e6b20717012a9ca49fa3bffb08a5db0bdd2e7}
\newcommand{\commitjapanese}{34fd0a9c71c6efdca2f154043820f2a6aef7fb3a}
\newcommand{\commitponceletthree}{9affe6203bbf0ec0f42eeb8c8ce567159de19b40}
\expandafter\def\csname commitof:carnot\endcsname{\commitcarnot}
\expandafter\def\csname commitof:viviani\endcsname{\commitviviani}
\expandafter\def\csname commitof:degua\endcsname{\commitdegua}
\expandafter\def\csname commitof:japanese\endcsname{\commitjapanese}
\expandafter\def\csname commitof:poncelet3\endcsname{\commitponceletthree}
\newcommand{\commitof}[1]{\csname commitof:#1\endcsname}
\expandafter\def\csname repoof:carnot\endcsname{carnot}
\expandafter\def\csname repoof:viviani\endcsname{viviani}
\expandafter\def\csname repoof:degua\endcsname{degua}
\expandafter\def\csname repoof:japanese\endcsname{japanese-theorem}
\expandafter\def\csname repoof:poncelet3\endcsname{poncelet-triangles}
\newcommand{\repo}[1]{https://github.com/jpbald93/\csname repoof:#1\endcsname}
\newcommand{\src}[4]{\href{\repo{#1}/blob/\commitof{#1}/#2/#3\#L#4}{\texttt{#3}:#4}}

\lstdefinelanguage{lean}{
  morekeywords={theorem,lemma,def,structure,inductive,where,by,fun,noncomputable},
  sensitive=true, morecomment=[l]{--}, morecomment=[s]{/-}{-/}}
\lstset{language=lean, basicstyle=\ttfamily\small, keywordstyle=\bfseries,
  columns=fullflexible, keepspaces=true, breaklines=true, frame=single, framerule=0.3pt,
  literate={ℝ}{{$\mathbb{R}$}}1 {ℕ}{{$\mathbb{N}$}}1 {≠}{{$\neq$}}1 {↔}{{$\leftrightarrow$}}2
    {→}{{$\to$}}1 {∀}{{$\forall$}}1 {∃}{{$\exists$}}1 {∧}{{$\wedge$}}1 {∈}{{$\in$}}1
    {∑}{{$\Sigma$}}1 {≤}{{$\le$}}1 {ᵥ}{{$_v$}}1 {Γ}{{$\Gamma$}}1 {ω}{{$\omega$}}1
    {∘}{{$\circ$}}1 {×}{{$\times$}}1}

\title{Five classical theorems on triangles and simplices,\\ formalised in Lean~4 with Mathlib:\\
Carnot, Viviani, De~Gua, the Japanese theorem and Poncelet's porism}
\author{Josh Bald\\ \small Independent researcher, Ontario, Canada\\ \small\texttt{jpbald93@gmail.com}}
\date{30 September 2026\\ \small Version~2, October 2026: adds Layers~14 and~15 (Section~\ref{sec:l1415})}

\emergencystretch=3em
\begin{document}
\maketitle

\begin{abstract}
We report Lean~4 formalisations, built on Mathlib, of five classical results of elementary
geometry: Carnot's theorem $\sum_i d_i(O)=R+r$ for a triangle; Viviani's theorem, generalised to
simplices of every positive dimension as the identity $\sum_i d_i(p)/h_i=1$; De~Gua's theorem in every
dimension, for simplices whose edges at one vertex are pairwise orthogonal (with $k$-volume defined by a Gram determinant), together with the projection form obtained from a
self-contained proof of the Cauchy--Binet formula; the Japanese theorem for convex cyclic polygons,
for all triangulations by maximal sets of non-crossing diagonals; and Poncelet's porism for
triangles, with its converse, in the ``construction'' form. The developments follow two strands.
Viviani, Carnot and the Japanese theorem are statements about signed distances to the facets of a
simplex, equivalently barycentric coordinates, and Poncelet's porism rests on the same incentre and
excentre vocabulary together with circle tangency; here Mathlib's \lean{Affine.Simplex} API
(\lean{signedInfDist}, \lean{height}, \lean{incenter}, \lean{excenter}, \lean{circumcenter})
supplies most of what is needed. De~Gua's theorem belongs to a second strand, Gram determinants of
edge vectors. We state every hypothesis exactly,
explain why each is present, and give lemmas that tie the formal definitions to the familiar
informal ones. Together the developments have about 3\,200 lines of new Lean code (plus about
900 lines of verbatim copies; the Layer-14/15 check files added in version~2 have about 2\,200 more), with no \code{sorry} and only Lean's standard axioms. As far as we
could determine, the geometric signed-distance formulation of Carnot's theorem, the
$n$-dimensional Viviani and De~Gua theorems in this form, the Japanese theorem and Poncelet's
porism have no earlier complete Lean formalisation; closely related Lean material does exist for
the trigonometric identity underlying Carnot's theorem, for the three-dimensional De~Gua theorem,
for coordinate versions of Viviani's theorem and for the Cauchy--Binet formula, and we describe it
in Section~\ref{sec:related}.
\end{abstract}

\section{Introduction}

The five theorems of this paper are classical, short, and familiar from olympiad geometry.
They are all listed in Mathlib's ``1000+ theorems'' file~\cite{Mathlib1000}, and at the Mathlib
version we use none of them has a declaration recorded there. They are:
\begin{itemize}
\item \emph{Carnot's theorem}. In a triangle with circumcentre $O$, circumradius $R$ and
inradius $r$, the signed distances from $O$ to the three side lines sum to $R+r$.
\item \emph{Viviani's theorem}. In an equilateral triangle the distances from an interior point to
the three sides sum to the altitude.
\item \emph{De~Gua's theorem}. If the three edges of a tetrahedron at one vertex are pairwise
perpendicular, the square of the area of the opposite face is the sum of the squares of the areas
of the other three faces~\cite{ConantBeyer}.
\item \emph{The Japanese theorem}. For a convex polygon inscribed in a circle, the sum of the
inradii of the triangles of a triangulation does not depend on the triangulation~\cite{FukagawaPedoe}.
\item \emph{Poncelet's porism for triangles}. Let $\Gamma$ and $\omega$ be circles with radii
$R>r>0$ and centres at distance $d$. If $d^2=R(R-2r)$, then every point of $\Gamma$ is a vertex of a
triangle inscribed in $\Gamma$ and circumscribed about $\omega$. Conversely, if the centre of
$\omega$ is inside~$\Gamma$ and one such triangle exists, then $d^2=R(R-2r)$~\cite{Flatto}.
In this paper ``circumscribed about $\omega$'' means that the three side \emph{lines} are tangent
to~$\omega$; see Section~\ref{sec:poncelet} for how this relates to the classical reading.
\end{itemize}

\paragraph{Two strands.}
Four of the five theorems use the vocabulary of \emph{signed distances} from a point to the facets
of a simplex; De~Gua's theorem uses Gram determinants instead. For a simplex $s$ with vertices $p_0,\dots,p_n$, let $d_i(p)$ be the signed distance
from $p$ to the hyperplane of the facet opposite $p_i$, positive on the side of $p_i$, and let
$h_i=d_i(p_i)$ be the $i$-th height. If $p=\sum_i w_ip_i$ with $\sum_iw_i=1$, then
\begin{equation}\label{eq:bary}
d_i(p)=w_i\,h_i .
\end{equation}
Mathlib provides $d_i$ as \lean{Affine.Simplex.signedInfDist} and \eqref{eq:bary} as
\lean{Affine.Simplex.signedInfDist_affineCombination}. The roles of the two strands are as follows:
\begin{itemize}
\item \emph{Viviani} is the statement $\sum_i d_i/h_i=\sum_iw_i=1$, specialised to equal heights.
\item \emph{Carnot} is \eqref{eq:bary} for the circumcentre, whose barycentric weights are known,
followed by a polynomial identity in the side lengths.
\item \emph{The Japanese theorem} is Carnot's theorem applied to each triangle of the
triangulation. The contributions of each diagonal cancel, since a diagonal is traversed in
opposite directions by its two triangles and the signed distance changes sign.
\item \emph{Poncelet's porism} for triangles does not follow from \eqref{eq:bary}: the forward
direction is an explicit coordinate construction of tangent lines to a circle. What it shares with
the strand above is the incentre/excentre vocabulary. Its converse rests on Euler's formula
$OI^2=R(R-2r)$, a statement about the incentre, the point with $d_0=d_1=d_2=r$: the centre of
$\omega$ is equidistant from the three side lines, so it is the incentre or an excentre.
\item \emph{De~Gua's theorem} forms a second strand. Its natural language is Gram determinants of
edge vectors rather than signed distances, and the general form is the Cauchy--Binet formula. It is
a companion result, not another application of~\eqref{eq:bary}.
\end{itemize}
The actual proof reuse is Carnot $\to$ Japanese (the Japanese development imports a copy of the
Carnot files), and Viviani and the Poncelet converse share Mathlib's simplex, incentre and
excentre infrastructure.

\paragraph{Contributions.}
\begin{itemize}
\item Lean formalisations of all five theorems, with the exact statements given in
Sections~\ref{sec:carnot}--\ref{sec:poncelet} and indexed in Table~\ref{tab:thms}.
\item Generalisations where they come cheaply: Viviani for simplices of every positive dimension, with a
characterisation of the simplices for which the distance sum is constant; De~Gua in every
dimension and in every real inner product space, with volume given by Definition~\ref{def:volsq}; the Japanese theorem for \emph{every}
triangulation, defined geometrically as a maximal set of chords with pairwise disjoint interiors;
Poncelet's porism in the form ``for $\omega$ centred inside $\Gamma$, the construction from one
point closes if and only if it closes from every point, if and only if $d^2=R(R-2r)$''.
\item A self-contained Lean proof of the Cauchy--Binet formula for rectangular matrices over a
commutative ring, which is not in Mathlib at the version we use (Section~\ref{sec:degua}).
\item Correctness lemmas that connect each formal definition with its informal reading
(Section~\ref{sec:correct}).
\end{itemize}

\paragraph{Availability.}
{\sloppy The five developments are separate Lake projects, published at
\url{\repo{carnot}}, \url{\repo{viviani}}, \url{\repo{degua}}, \url{\repo{japanese}} and
\url{\repo{poncelet3}}. Line references in this paper are to commits \texttt{98e5614} (Carnot),
\texttt{029aa3d} (Viviani), \texttt{d46e6b2} (De~Gua), \texttt{34fd0a9} (Japanese) and
\texttt{9affe62} (Poncelet). Each repository has a README with build and check commands.
All five use Lean \texttt{v4.33.1} and the same Mathlib revision
(\texttt{0df444a3}, tag \texttt{v4.33.1}), pinned in \texttt{lake-manifest.json}.
The Japanese development contains a verbatim copy of the Carnot development under
\texttt{Japanese/Carnot/} (only the import prefix differs). The Poncelet development contains, under
\texttt{Poncelet3/Euler/}, a verbatim copy of three files from the author's separate formalisation
of Feuerbach's theorem, which provides Euler's formula (again only the import line differs).
Version~2 of this paper adds the Layer-14 and Layer-15 files of Section~\ref{sec:l1415} in a new
commit of each repository: \texttt{fa8c37951d290d26db66ff83c358b09b156d26e6} (Carnot), \texttt{e71460f2edb224b2e7e98a39c6db6bead93613cc} (Viviani),
\texttt{daa89522f4a85e6afb8fe4d4d1a94b3d0761943e} (De~Gua), \texttt{be4503b4e1dfdf7e527e2425276f992ecabb4bc2} (Japanese) and \texttt{aa9da6c4e077cee2e9191ad9011a246ad9e04f2f} (Poncelet). The
earlier statements and library files are unchanged in that commit (only the root file gains two
import lines, and \path{gate.sh} and the README are updated), so the line references above are valid
at both commits. Supporting material (the Layer-14 and Layer-15 reports, the blind writer's statements and report, the transcript audit, and the reviews of this version with the fixes made in response) is in the folder \texttt{Triangle-Simplex} of the repository \url{https://github.com/jpbald93/LEAN-verification}.\par}

{\small\begin{longtable}{@{}>{\raggedright\arraybackslash\hangindent=1.5em\hangafter=1}p{0.52\textwidth}p{0.16\textwidth}>{\raggedright\arraybackslash}p{0.26\textwidth}@{}}
\caption{Main formal statements, grouped by repository. Links point to the declaration line
at the commit of that repository listed under Availability.}\label{tab:thms}\\
\toprule
Lean name & In this paper & Source \\ \midrule
\endfirsthead
\multicolumn{3}{@{}l}{\small\emph{Table~\ref{tab:thms}, continued}}\\
\toprule
Lean name & In this paper & Source \\ \midrule
\endhead
\bottomrule
\endfoot
\multicolumn{3}{@{}l}{\emph{Repo \texttt{carnot}}}\\
\lean{Carnot.signedSideDist} & Def.~\ref{def:ssd} & \src{carnot}{Carnot}{Main.lean}{27} \\
\lean{Carnot.signedSideDist_eq_signedInfDist} & Lemma~\ref{lem:carnot-correct} & \src{carnot}{Carnot}{Main.lean}{49} \\
\lean{Carnot.signedSideDist_incenter} & Lemma~\ref{lem:carnot-correct} & \src{carnot}{Carnot}{Main.lean}{65} \\
\lean{Carnot.carnot} & Theorem~\ref{thm:carnot} & \src{carnot}{Carnot}{Main.lean}{75} \\
\lean{Carnot.carnot_sum} & Theorem~\ref{thm:carnot} & \src{carnot}{Carnot}{Main.lean}{97} \\
\lean{Carnot.signedSideDist_circumcenter_eq} & Theorem~\ref{thm:carnot-cos} & \src{carnot}{Carnot}{Main.lean}{113} \\
\lean{Carnot.cos_add_cos_add_cos} & Theorem~\ref{thm:carnot-cos} & \src{carnot}{Carnot}{Main.lean}{151} \\
\lean{Carnot.carnot_poly} & Section~\ref{sec:carnot} & \src{carnot}{Carnot}{Algebra.lean}{32} \\
\multicolumn{3}{@{}l}{\emph{Repo \texttt{viviani}}}\\
\lean{Affine.Simplex.sum_signedInfDist_div_height} & Theorem~\ref{thm:viv-weights} & \src{viviani}{Viviani}{Basic.lean}{37} \\
\lean{Affine.Simplex.sum_signedInfDist_eq_height} & Theorem~\ref{thm:viv-heights} & \src{viviani}{Viviani}{Heights.lean}{34} \\
\lean{Affine.Simplex.exists_forall_sum_signedInfDist_eq_iff} & Theorem~\ref{thm:viv-heights} & \src{viviani}{Viviani}{Heights.lean}{43} \\
\lean{Affine.Simplex.signedInfDist_eq_dist_of_mem_closedInterior} & Lemma~\ref{lem:viv-unsigned} & \src{viviani}{Viviani}{Heights.lean}{64} \\
\lean{Affine.Simplex.sum_dist_eq_height} & Theorem~\ref{thm:viv-heights} & \src{viviani}{Viviani}{Heights.lean}{73} \\
\lean{Affine.Simplex.Equilateral.height_eq} & Proposition~\ref{prop:equi} & \src{viviani}{Viviani}{Equilateral.lean}{62} \\
\lean{Affine.Simplex.Regular.sum_dist_eq_height} & Corollary~\ref{cor:viviani} & \src{viviani}{Viviani}{Equilateral.lean}{92} \\
\lean{Viviani.equilateral_triangle'} & Corollary~\ref{cor:viviani} & \src{viviani}{Viviani}{Equilateral.lean}{109} \\
\lean{Affine.Simplex.Equilateral.height_eq} + \lean{Affine.Simplex.sum_dist_eq_height} & Corollary~\ref{cor:viviani} (equilateral case) & \src{viviani}{Viviani}{Equilateral.lean}{62}, \src{viviani}{Viviani}{Heights.lean}{73} \\
\multicolumn{3}{@{}l}{\emph{Repo \texttt{degua}}}\\
\lean{DeGua.de_gua_three} & Theorem~\ref{thm:degua3} & \src{degua}{DeGua}{Three.lean}{42} \\
\lean{DeGua.simplexVolSq} & Def.~\ref{def:volsq} & \src{degua}{DeGua}{General.lean}{97} \\
\lean{DeGua.det_diagonal_add_const} & Lemma~\ref{lem:mdl} & \src{degua}{DeGua}{General.lean}{45} \\
\lean{DeGua.de_gua} & Theorem~\ref{thm:degua} & \src{degua}{DeGua}{General.lean}{130} \\
\lean{DeGua.triAreaSq_eq_simplexVolSq} & Lemma~\ref{lem:degua-bridge} & \src{degua}{DeGua}{Bridge.lean}{16} \\
\lean{DeGua.de_gua_three_of_general} & Lemma~\ref{lem:degua-bridge} & \src{degua}{DeGua}{Bridge.lean}{28} \\
\lean{DeGua.det_mul_eq_sum_strictMono} & Theorem~\ref{thm:cb} & \src{degua}{DeGua}{CauchyBinet.lean}{33} \\
\lean{DeGua.det_gram_eq_sum_sq_minors} & Corollary~\ref{cor:proj} & \src{degua}{DeGua}{CauchyBinet.lean}{91} \\
\multicolumn{3}{@{}l}{\emph{Repo \texttt{japanese}}}\\
\lean{Japanese.chordDist} & Def.~\ref{def:chord} & \src{japanese}{Japanese}{Basic.lean}{30} \\
\lean{Japanese.chordDist_swap} & Lemma~\ref{lem:swap} & \src{japanese}{Japanese}{Basic.lean}{35} \\
\lean{Japanese.triInradius_eq} & Lemma~\ref{lem:carnot-chord} & \src{japanese}{Japanese}{Basic.lean}{130} \\
\lean{Japanese.japanese_quadrilateral} & Theorem~\ref{thm:jq} & \src{japanese}{Japanese}{Basic.lean}{193} \\
\lean{Japanese.Triangulation} & Def.~\ref{def:tree} & \src{japanese}{Japanese}{Polygon.lean}{33} \\
\lean{Japanese.inradiusSum_eq} & Theorem~\ref{thm:jtree} & \src{japanese}{Japanese}{Polygon.lean}{97} \\
\lean{Japanese.japanese_polygon} & Theorem~\ref{thm:jtree} & \src{japanese}{Japanese}{Polygon.lean}{131} \\
\lean{Japanese.IsTriangulation} & Def.~\ref{def:triang} & \src{japanese}{Japanese}{Combinatorial.lean}{30} \\
\lean{Japanese.isTriangulation_iff_exists_tree} & Proposition~\ref{prop:bij} & \src{japanese}{Japanese}{Combinatorial.lean}{415} \\
\lean{Japanese.Triangulation.chords_injective} & Proposition~\ref{prop:bij} & \src{japanese}{Japanese}{Combinatorial.lean}{453} \\
\lean{Japanese.treeEquiv} & Proposition~\ref{prop:bij} & \src{japanese}{Japanese}{Combinatorial.lean}{481} \\
\lean{Japanese.toFinset_triangles} & Proposition~\ref{prop:bij} (triangles $=$ \lean{triangleSet}) & \src{japanese}{Japanese}{Combinatorial.lean}{259} \\
\lean{Japanese.card_triangleSet} & Proposition~\ref{prop:bij} ($j-i-1$ triangles) & \src{japanese}{Japanese}{Combinatorial.lean}{265} \\
\lean{Japanese.japanese_polygon_combinatorial} & Theorem~\ref{thm:japanese} & \src{japanese}{Japanese}{Combinatorial.lean}{520} \\
\lean{Japanese.sbtw_inter_iff_interleave} & Proposition~\ref{prop:cross} & \src{japanese}{Japanese}{Geometric.lean}{182} \\
\lean{Japanese.geomCrosses_iff} & Proposition~\ref{prop:cross} & \src{japanese}{Japanese}{Geometric.lean}{220} \\
\lean{Japanese.IsGeomTriangulation} & Def.~\ref{def:triang} & \src{japanese}{Japanese}{Geometric.lean}{209} \\
\lean{Japanese.japanese_polygon_geometric} & Theorem~\ref{thm:japanese} & \src{japanese}{Japanese}{Geometric.lean}{276} \\
\multicolumn{3}{@{}l}{\emph{Repo \texttt{poncelet3}}}\\
\lean{Feuerbach.dist_circumcenter_incenter_sq} & Theorem~\ref{thm:euler} & \src{poncelet3}{Poncelet3/Euler}{Main.lean}{100} \\
\lean{Feuerbach.dist_circumcenter_excenter_sq} & Theorem~\ref{thm:euler} & \src{poncelet3}{Poncelet3/Euler}{Main.lean}{105} \\
\lean{Poncelet3.exists_poncelet_triangle} & Theorem~\ref{thm:porism} & \src{poncelet3}{Poncelet3}{Porism.lean}{104} \\
\lean{Poncelet3.euler_of_isTangent} & Theorem~\ref{thm:converse} & \src{poncelet3}{Poncelet3}{Closure.lean}{35} \\
\lean{Poncelet3.poncelet_triangle_iff} & Theorem~\ref{thm:converse} & \src{poncelet3}{Poncelet3}{Closure.lean}{105} \\
\lean{Poncelet3.radius_lt_dist_center_of_mem} & Lemma~\ref{lem:pon-correct} & \src{poncelet3}{Poncelet3}{Construction.lean}{77} \\
\lean{Poncelet3.tangentsFrom_eq_pair} & Lemma~\ref{lem:pon-correct} & \src{poncelet3}{Poncelet3}{Construction.lean}{263} \\
\lean{Poncelet3.secondInter_spec} & Lemma~\ref{lem:pon-correct} & \src{poncelet3}{Poncelet3}{Construction.lean}{313} \\
\lean{Poncelet3.poncelet_construction} & Theorem~\ref{thm:construction} & \src{poncelet3}{Poncelet3}{Construction.lean}{334} \\
\lean{Poncelet3.ConstructionCloses} & Def.~\ref{def:closes} & \src{poncelet3}{Poncelet3}{Package.lean}{37} \\
\lean{Poncelet3.poncelet_construction_iff} & Theorem~\ref{thm:construction} & \src{poncelet3}{Poncelet3}{Package.lean}{98} \\
\lean{Poncelet3.poncelet_all_or_nothing} & Theorem~\ref{thm:construction} & \src{poncelet3}{Poncelet3}{Package.lean}{107} \\
\end{longtable}}

\section{Setting and Mathlib vocabulary}\label{sec:setting}

Throughout, $V$ is a real inner product space and $P$ a metric affine space over it:
\begin{center}\small
\code{[NormedAddCommGroup V] [InnerProductSpace $\mathbb{R}$ V] [MetricSpace P] [NormedAddTorsor V P]}
\end{center}
No finite-dimensionality is assumed unless stated. The Japanese theorem and Poncelet's porism
additionally assume \code{[Fact (finrank $\mathbb{R}$ V = 2)]}, so that $P$ is a Euclidean plane; the Japanese
theorem also fixes an orientation \code{o : Orientation $\mathbb{R}$ V (Fin 2)}, with area form
$\omega_o$ (\lean{Orientation.areaForm}).

We use the following Mathlib notions for a simplex \code{s : Affine.Simplex $\mathbb{R}$ P n}; an
\lean{Affine.Triangle} is the case $n=2$. Heights and signed facet distances require $n\ge1$
(Mathlib's \code{[NeZero n]}), since a $0$-simplex has no facets.
\begin{itemize}
\item \lean{Affine.Simplex.height} \code{s i} is the distance from the $i$-th vertex to the
affine span of the opposite facet; \lean{Affine.Simplex.height_pos} says it is positive.
\item \lean{Affine.Simplex.signedInfDist} \code{s i p} is the signed distance from $p$ to the
affine span of the facet opposite vertex $i$, positive towards vertex~$i$. Components of $p$
orthogonal to the span of the simplex are ignored.
\item \lean{Affine.Simplex.circumcenter}, \lean{Affine.Simplex.circumradius},
\lean{Affine.Simplex.incenter}, \lean{Affine.Simplex.inradius},
\lean{Affine.Simplex.excenter} and \lean{Affine.Simplex.exradius} have their usual meaning. The
incentre and excentres are defined for general simplices through weights
(\lean{Affine.Simplex.excenterWeights}).
\item \lean{Affine.Simplex.closedInterior} is the set of affine combinations with weights in
$[0,1]$, that is, the closed simplex.
\item \lean{Affine.Simplex.Equilateral} means all edges have the same length, and
\lean{Affine.Simplex.Regular} means every permutation of the vertices is induced by an isometry.
\end{itemize}
Circles are \lean{EuclideanGeometry.Sphere} \code{P}. An affine subspace is tangent to a sphere
(\lean{EuclideanGeometry.Sphere.IsTangent}) if it lies in the tangent space
(\lean{EuclideanGeometry.Sphere.orthRadius}) at some point of the sphere that it contains.
\lean{EuclideanGeometry.Sphere.tangentsFrom} \code{A} is the set of maximal tangent spaces
through $A$, and \lean{EuclideanGeometry.Sphere.secondInter} is the second intersection of a line
with a sphere. We write \code{line[$\mathbb{R}$, A, B]} for the affine span of $\{A,B\}$.

\section{Carnot's theorem}\label{sec:carnot}

\begin{definition}\label{def:ssd}
For a triangle $t$ and $i\in\{0,1,2\}$, let $F_i$ be the foot of the altitude from vertex $A_i$
and $h_i$ the height. The signed side distance of $p$ is
\[
\lean{signedSideDist}\ t\ i\ p=\langle p-F_i,\ A_i-F_i\rangle/h_i ,
\]
the component of $p-F_i$ along the unit normal to side~$i$ that points towards $A_i$.
\end{definition}

It agrees with Mathlib's \lean{signedInfDist} on the plane of the triangle
(Lemma~\ref{lem:carnot-correct}), which is where the circumcentre lies. The key property is \lean{Carnot.signedSideDist_affineCombination}, the
triangle case of~\eqref{eq:bary}.

\begin{theorem}[Carnot]\label{thm:carnot}
For every triangle $t$ in $P$,
\[
\sum_{i=0}^{2}\lean{signedSideDist}\ t\ i\ O=R+r,
\]
where $O$, $R$, $r$ are the circumcentre, circumradius and inradius of~$t$.
\end{theorem}

\begin{lstlisting}
theorem carnot :
    signedSideDist t 0 t.circumcenter + signedSideDist t 1 t.circumcenter +
        signedSideDist t 2 t.circumcenter = t.circumradius + t.inradius
\end{lstlisting}

There are no hypotheses beyond $t$ being a triangle, which in Mathlib includes affine
independence. There is no dimension hypothesis. The distances are signed, and this is essential:
for an obtuse triangle the circumcentre lies outside, and the distance to the side opposite the
obtuse angle is negative. Theorem~\ref{thm:carnot-cos} makes the sign convention visible.

\begin{theorem}\label{thm:carnot-cos}
Let $\alpha_i$ be the angle of $t$ at vertex $A_i$. Then
$\lean{signedSideDist}\ t\ i\ O=R\cos\alpha_i$ for each $i$, and
$\cos\alpha_0+\cos\alpha_1+\cos\alpha_2=1+r/R$.
\end{theorem}

\paragraph{Proof.}
The circumcentre is written as $O=\sum_iw_iA_i$. With side lengths $a,b,c$ and $k=2\cdot$area,
one has $h_i=k/a_i$, $2kR=abc$ and, from \lean{Affine.Simplex.inradius_eq_abs_inv_sum},
$r=k/(a+b+c)$. By \eqref{eq:bary}, Theorem~\ref{thm:carnot} becomes
$\sum_iw_ik/a_i=R+k/(a+b+c)$, which, after clearing denominators, is the polynomial identity
\lean{Carnot.carnot_poly}:
\[
\bigl(w_0bc+w_1ca+w_2ab\bigr)(a+b+c)=2R^2(a+b+c)+abc .
\]
It holds modulo the three equations $|O-A_i|^2=R^2$, written in the side lengths and the weights.
The coefficients of the \code{linear\_combination} certificate (of degree two) were found with a
computer algebra system. The cosine form uses three further certificates of degree three and the
law of cosines from Mathlib. The frame machinery
(\lean{Carnot.frame_data}, \lean{Carnot.height_sq_mul_sq}) is adapted from the author's
Feuerbach development.

\section{Viviani's theorem, generalised to simplices}\label{sec:viviani}

Here $s$ is an $n$-simplex with $n\ge1$ (\code{[NeZero n]}), in any real inner product space. We
write $d_i(p)$ for \code{s.signedInfDist i p} and $h_i$ for \code{s.height i}.

\begin{theorem}\label{thm:viv-weights}
For every point $p$ in the affine span of $s$, $\sum_i d_i(p)/h_i=1$.
Moreover, $d_i(p)/h_i$ is the $i$-th barycentric coordinate of~$p$.
\end{theorem}

\begin{proof}
Write $p=\sum_iw_ip_i$ with $\sum_iw_i=1$, and apply \eqref{eq:bary}.
\end{proof}

\begin{theorem}\label{thm:viv-heights}
Suppose all heights of $s$ are equal to $h$. Then $\sum_id_i(p)=h$ for every $p$ in the affine
span of~$s$, and $\sum_i\operatorname{dist}(p,\Pi_i)=h$ for every $p$ in the closed simplex, where
$\Pi_i$ is the affine span of the $i$-th facet. Conversely, if $p\mapsto\sum_id_i(p)$ is constant on
the affine span of $s$, then all heights are equal.
\end{theorem}

\begin{lstlisting}
theorem sum_signedInfDist_eq_height (hh : ∀ i j, s.height i = s.height j)
    {p : P} (hp : p ∈ affineSpan ℝ (Set.range s.points)) (k : Fin (n + 1)) :
    ∑ i, s.signedInfDist i p = s.height k
theorem exists_forall_sum_signedInfDist_eq_iff :
    (∃ c : ℝ, ∀ p ∈ affineSpan ℝ (Set.range s.points), ∑ i, s.signedInfDist i p = c) ↔
      ∀ i j, s.height i = s.height j
\end{lstlisting}

The converse is proved by evaluating the sum at two vertices: $\sum_id_i(p_j)=h_j$.

\paragraph{Why each hypothesis.}
The signed form holds on the whole affine span, which is the natural domain of~$d_i$. The unsigned
form needs $p$ in the closed simplex, where all $d_i(p)\ge0$ (weights in $[0,1]$); outside it, some
distance must be counted negatively. We use Mathlib's \lean{Affine.Simplex.closedInterior}, so the
unsigned form includes boundary points. ``Distance to a side'' means the distance to the affine span of the
facet (the side \emph{line} of a triangle), attained at the orthogonal projection
(\lean{Affine.Simplex.orthogonalProjectionSpan}). This is the classical reading. For an
equilateral triangle the foot of the perpendicular from a point of the closed triangle lies on the
side itself, but for a general simplex it need not, and we do not use or prove any such fact.

\begin{proposition}\label{prop:equi}
All heights of an equilateral simplex are equal.
\end{proposition}

\begin{proof}
For $i\ne j$, reflect in the perpendicular bisector hyperplane of $[p_i,p_j]$. This affine isometry
(\lean{Viviani.perpReflection}) swaps $p_i$ and $p_j$ and fixes every other vertex, since those are
equidistant from $p_i$ and $p_j$. Heights are invariant under isometries
(\lean{Affine.Simplex.height_map}) and permute under reindexing
(\lean{Affine.Simplex.height_reindex}), so $h_i=h_j$.
\end{proof}

\begin{corollary}[Viviani]\label{cor:viviani}
For a regular (or just equilateral) simplex, and every point $p$ of the closed simplex, the
distances from $p$ to the facet hyperplanes sum to the height. For an equilateral triangle
\code{t} and \code{p $\in$ t.closedInterior}, the three distances sum to \code{t.height 0}.
\end{corollary}

Equal heights do not imply equilateral when $n\ge3$. Since $h_i\cdot\operatorname{vol}(F_i)=
n\operatorname{vol}(s)$, equal heights means equal facet volumes. In dimension three this holds for
every tetrahedron with four congruent faces (a disphenoid), which need not be regular. So
Theorem~\ref{thm:viv-heights} characterises the simplices with constant distance sum: they are the
simplices with equal heights, a strictly larger class than the equilateral ones. The identity
$h_i\operatorname{vol}(F_i)=n\operatorname{vol}(s)$ is not formalised here.

\section{De~Gua's theorem and the Cauchy--Binet formula}\label{sec:degua}

\subsection{The three-dimensional statement}
In $\RR^3$ (\code{EuclideanSpace $\mathbb{R}$ (Fin 3)}), let
\lean{DeGua.triAreaSq}~$X\,Y\,Z=|(Y-X)\times(Z-X)|^2/4$, the squared area of a triangle.

\begin{theorem}\label{thm:degua3}
If $A-O$, $B-O$, $C-O$ are pairwise orthogonal, then
$\lean{triAreaSq}\,A\,B\,C=\lean{triAreaSq}\,O\,A\,B+\lean{triAreaSq}\,O\,B\,C+\lean{triAreaSq}\,O\,C\,A$.
\end{theorem}

The proof is one \code{linear\_combination} (\lean{DeGua.cross_sub_sq}).

\subsection{All dimensions}
There is no cross product in general, so we measure volume by Gram determinants.

\begin{definition}\label{def:volsq}
For points $x_0,\dots,x_k$ of a real inner product space $E$,
\[
\lean{simplexVolSq}\ x=\det\bigl(\langle x_{i}-x_0,\,x_{j}-x_0\rangle\bigr)_{1\le i,j\le k}\big/(k!)^2 .
\]
\end{definition}

This is the standard squared $k$-volume of the simplex $x_0\cdots x_k$. We stress that it is a
\emph{definition}: we have not related it to a measure. Mathlib now has a Euclidean Hausdorff measure
(\lean{MeasureTheory.Measure.euclideanHausdorffMeasure}), and connecting the two would need the
facts that a parallelotope has volume $\sqrt{\det G}$ and a simplex $1/k!$ of that.

\begin{theorem}[De~Gua, general form]\label{thm:degua}
Let $O\in E$ and $P_0,\dots,P_m\in E$ with $P_i-O$ pairwise orthogonal. Then
\[
\lean{simplexVolSq}\,(P_0,\dots,P_m)=\sum_{i=0}^m\lean{simplexVolSq}\,(O,P_0,\dots,\widehat{P_i},\dots,P_m).
\]
\end{theorem}

\begin{lstlisting}
theorem de_gua (O : E) {m : ℕ} (P : Fin (m + 1) → E)
    (h : Pairwise fun i j => inner ℝ (P i - O) (P j - O) = 0) :
    simplexVolSq P =
      ∑ i : Fin (m + 1), simplexVolSq (Fin.cons O (P ∘ i.succAbove) : Fin (m + 1) → E)
\end{lstlisting}

The only hypothesis is orthogonality. The points need not be distinct and $E$ need not be
finite-dimensional, and degenerate configurations (for instance a zero edge) are allowed.

\begin{lemma}\label{lem:mdl}
For a finite index type $\iota$, $b\in\RR^\iota$ and $c\in\RR$, with $J$ the all-ones matrix,
$\det(\operatorname{diag}(b)+cJ)=\prod_ib_i+c\sum_i\prod_{j\ne i}b_j$.
\end{lemma}

When all $b_i\ne0$, this follows from Mathlib's matrix determinant lemma
\lean{Matrix.det_one_add_replicateCol_mul_replicateRow}. For general $b$, replace $b$ by
$b+t\mathbf 1$. Both sides are continuous in $t$, and they agree whenever $t\ne-b_i$ for all $i$,
that is, on the complement of a finite set. That complement is dense, so they agree at $t=0$. The
Gram matrix of the edges $P_j-P_0$ of the facet opposite $O$ is
$\operatorname{diag}(|P_j-O|^2)_{j\ge1}+|P_0-O|^2J$,
so Lemma~\ref{lem:mdl} gives $\sum_i\prod_{k\ne i}|P_k-O|^2$. The facets through $O$ have diagonal
Gram matrices with determinants $\prod_{k\ne i}|P_k-O|^2$.

\subsection{Cauchy--Binet and the projection form}

\begin{theorem}[Cauchy--Binet]\label{thm:cb}
Let $R$ be a commutative ring, $A$ a $k\times n$ and $B$ an $n\times k$ matrix over $R$. Then
\[
\det(AB)=\sum_{S}\det(A_{\cdot,S})\det(B_{S,\cdot}),
\]
the sum over strictly increasing $S:\{0..k-1\}\to\{0..n-1\}$, i.e.\ over $k$-subsets of columns.
\end{theorem}

\begin{lstlisting}
theorem det_mul_eq_sum_strictMono (A : Matrix (Fin k) (Fin n) R)
    (B : Matrix (Fin n) (Fin k) R) :
    det (A * B) = ∑ s ∈ univ.filter (fun s : Fin k → Fin n => StrictMono s),
      det (A.submatrix id s) * det (B.submatrix s id)
\end{lstlisting}

The proof expands $\det(AB)$ as a sum over all maps $p:\mathrm{Fin}\,k\to\mathrm{Fin}\,n$, drops
the non-injective ones (their minor has two equal columns), and re-indexes the injective ones as
pairs (strictly monotone $s$, permutation $\tau$) with Mathlib's \lean{Tuple.sort} and
\lean{Tuple.unique_monotone}. It covers $k>n$ (both sides vanish) and $k=0$.

\begin{corollary}[Pythagorean theorem for $k$-volumes]\label{cor:proj}
For $v_0,\dots,v_{k-1}\in\RR^n$, the Gram determinant of $v$ equals the sum, over all coordinate
$k$-planes, of the squared $k\times k$ minors of the matrix with rows $v_i$. That is, the squared
volume of the parallelotope spanned by $v$ is the sum of the squared volumes of its orthogonal
projections onto the coordinate $k$-planes.
\end{corollary}

For $k=n-1$, applied to the facet opposite the corner of a simplex whose corner edges lie along
the coordinate axes, this recovers De~Gua's theorem; this is Conant and Beyer's
generalisation~\cite{ConantBeyer}. We have not formalised that deduction:
Theorem~\ref{thm:degua} is proved directly.
Theorem~\ref{thm:cb} is, as far as we could determine, absent from Mathlib at our pinned version,
and a Mathlib pull request for it is open (Section~\ref{sec:related}). Ours may be useful as a
second, short proof ($\sim$70 lines); we make no stronger claim.

\section{The Japanese theorem}\label{sec:japanese}

Here $P$ is an oriented Euclidean plane.

\begin{definition}\label{def:chord}
For points $O,p,q$, let $\lean{chordDist}\ O\ p\ q=\omega_o(q-p,\,O-p)/|pq|$, for $p\ne q$ the signed distance
from $O$ to the directed line $pq$, positive if $O$ is to the left of the chord $p\to q$ (the total Lean definition returns $0$ when $p=q$). Let
\lean{Japanese.triInradius}~$a\,b\,c$ be the inradius of the triangle $abc$, or $0$ if $a,b,c$ are
collinear.
\end{definition}

\begin{lemma}[cancellation]\label{lem:swap}
$\lean{chordDist}\ O\ q\ p=-\lean{chordDist}\ O\ p\ q$.
\end{lemma}

\begin{lemma}[Carnot, chord form]\label{lem:carnot-chord}
If $a,b,c$ lie on a circle with centre $O$ and radius $R$ and $\omega_o(b-a,c-a)>0$, then
$\lean{triInradius}\ a\ b\ c=d(ab)+d(bc)+d(ca)-R$, where $d(pq)=\lean{chordDist}\ O\ p\ q$.
\end{lemma}

This is Theorem~\ref{thm:carnot} (from the copied Carnot files), after identifying the circle with
the circumcircle (the triangle spans the plane) and Carnot's inward signed distance with the chord
distance along the counterclockwise side (\lean{Japanese.signedSideDist_eq_chordDist}).

\begin{theorem}[quadrilateral]\label{thm:jq}
If $A,B,C,D$ lie on a circle and the triangles $ABC$, $ACD$, $ABD$, $BCD$ are positively oriented,
then $r(ABC)+r(ACD)=r(ABD)+r(BCD)$.
\end{theorem}

\paragraph{Polygons and triangulations.}
A polygon is a map $p:\mathbb N\to P$ of which only $p_0,\dots,p_{n-1}$ are used. The hypotheses are
\begin{itemize}
\item \code{hs : $\forall$ l < n, p l $\in$ s} for a circle \code{s}, and
\item \code{hconv}: $\omega_o(p_b-p_a,\ p_c-p_a)>0$ whenever $a<b<c<n$.
\end{itemize}
For $n\ge3$ the second says that the polygon is strictly convex and listed counterclockwise. For
$n=2$ it is vacuous (even $p_0=p_1$ is allowed); the formal theorems handle this harmless
degenerate case through empty sums of triangles. The hypothesis is needed because
Lemma~\ref{lem:carnot-chord} requires each triangle to be positively oriented. For a negatively
oriented triangle the three chord distances change sign, so the chord sum is $-(R+r)$ and the
formula becomes $r=-d(ab)-d(bc)-d(ca)-R$; reversing the ambient orientation restores the stated
form. The counterclockwise choice is thus a convention needed by the displayed signed formula, not
a geometric restriction: the inradius sum itself does not depend on the orientation, and the
clockwise case follows by reversing~$o$. For points on a circle, \code{hconv} is the same as
being in counterclockwise cyclic order, but we did not formalise this equivalence.

\begin{definition}\label{def:tree}
\lean{Japanese.Triangulation}~$i\ j$ is the inductive type with constructors
\code{edge i : Triangulation i (i+1)} and
\code{node (i<k) (k<j) : Triangulation i k $\to$ Triangulation k j $\to$ Triangulation i j}.
Its triangles (\lean{Japanese.Triangulation.triangles}) are the triples $(i,k,j)$ at the nodes.
\end{definition}

A \code{node} is the triangle $(i,k,j)$ standing on the side $ij$, together with triangulations
of the two remaining sub-polygons. Summing Lemma~\ref{lem:carnot-chord} over the triangles, the
diagonals $ik$ and $kj$ cancel by Lemma~\ref{lem:swap}. Induction on the tree gives:

\begin{theorem}\label{thm:jtree}
Under \code{hs} and \code{hconv}, every $T:\lean{Triangulation}\ 0\ (n-1)$ has inradius sum
$\sum_{l}d(p_lp_{l+1})-(n-2)R$, where the sum runs over the $n$ sides of the polygon. In particular any
two such $T,T'$ have the same inradius sum.
\end{theorem}

No flip-connectivity argument is needed: the value is computed for every tree directly.

\paragraph{Triangulations as sets of chords.}
The tree type is a convenient model, but a triangulation is usually defined as a maximal set of
non-crossing diagonals. We prove that the two agree. A chord of $i..j$ is a pair $(a,b)$ with
$i\le a<b\le j$; the chords $(a,b)$ and $(c,d)$ \emph{cross} if $a<c<b<d$ or $c<a<d<b$
(\lean{Japanese.Crosses}).

\begin{definition}\label{def:triang}
$E$ is a \emph{triangulation} of $i..j$ (\lean{Japanese.IsTriangulation}) if all elements are
chords, no two cross, and every chord that crosses no element of $E$ is in $E$. It is a
\emph{geometric triangulation} (\lean{Japanese.IsGeomTriangulation}) if the same holds with
``cross'' replaced by ``the open segments $p_ap_b$ and $p_cp_d$ meet'' for $e\ne f$. The triangles
of $E$ (\lean{Japanese.triangleSet}) are the triples $(a,b,c)$ with $ab$, $bc$, $ac\in E$.
\end{definition}

Sides of the polygon and the root side cross nothing, so by maximality they belong to~$E$. Hence
$E$ minus the sides is exactly a maximal non-crossing set of diagonals, and triangles can be read
off uniformly. Deriving the triangles from $E$, instead of axiomatising them, avoids conditions such
as ``each diagonal lies in exactly two triangles''.

\begin{proposition}\label{prop:bij}
For $i<j$, the map $T\mapsto\lean{Japanese.Triangulation.chords}\ T$ is a bijection from trees to
triangulations of $i..j$ (\lean{Japanese.treeEquiv}). Under it, the tree's triangles are exactly
\lean{triangleSet}, and every triangulation of $i..j$ has $j-i-1$ triangles.
\end{proposition}

The hard direction builds a tree from $E$. If $i+1<j$, let $k$ be the largest $b<j$ with
$(i,b)\in E$. Then $(k,j)$ crosses no element of $E$, so it is in $E$, and every chord of $E$ other
than the root side $(i,j)$ lies on one side of $k$ (\lean{Japanese.IsTriangulation.exists_apex}).
Restricting $E$ to both sides gives triangulations of $i..k$ and $k..j$. Strong
induction on $j-i$ finishes the proof.

\begin{proposition}\label{prop:cross}
Under \code{hconv}, for distinct chords $e,f$ of $0..n-1$, the open segments of $e$ and $f$ meet if and
only if $e$ and $f$ cross. Consequently, geometric and combinatorial triangulations coincide.
\end{proposition}

For interleaved indices the intersection point is written down explicitly as a ratio of area forms;
for non-interleaved disjoint pairs both endpoints of one chord lie strictly on one side of the
other; chords with a shared endpoint never meet in their interiors because three vertices are never
collinear. The case $e=f$ is excluded because a segment meets itself.

\begin{theorem}[Japanese theorem]\label{thm:japanese}
Let $n\ge2$, and assume \code{hs} and \code{hconv}. For any two geometric triangulations $E,E'$ of the
polygon $p_0\cdots p_{n-1}$,
$\sum_{t\in\lean{triangleSet}\,E}r(t)=\sum_{t\in\lean{triangleSet}\,E'}r(t)$. The same holds for
combinatorial triangulations.
\end{theorem}

\begin{lstlisting}
theorem japanese_polygon_geometric (s : Sphere P) (hn : 2 ≤ n)
    (hs : ∀ l < n, p l ∈ s)
    (hconv : ∀ a b c, a < b → b < c → c < n →
      0 < o.areaForm (p b -ᵥ p a) (p c -ᵥ p a))
    {E E' : Finset (ℕ × ℕ)} (hE : IsGeomTriangulation p 0 (n - 1) E)
    (hE' : IsGeomTriangulation p 0 (n - 1) E') :
    triangleInradiusSum p (triangleSet E) = triangleInradiusSum p (triangleSet E')
\end{lstlisting}

The hypothesis $n\ge2$ only excludes degenerate cases. What is still not formalised is a
dissection in the sense of ``a finite family of triangles with disjoint interiors whose union is the
polygon''; our geometric triangulations are maximal sets of non-crossing chords
(Section~\ref{sec:limits}).

\section{Poncelet's porism for triangles}\label{sec:poncelet}

Here $P$ is a Euclidean plane, $\Gamma$ and $\omega$ are spheres with centres $O$, $I$ and radii
$R$, $r$, and $d=|OI|$.

\begin{theorem}[Euler]\label{thm:euler}
For every triangle, $|OI|^2=R(R-2r)$ and $|OI_i|^2=R(R+2r_i)$, where $I_i$ and $r_i$ are the
excentre and exradius opposite vertex~$i$.
\end{theorem}

The two formulas are
\lean{Feuerbach.dist_circumcenter_incenter_sq} and \lean{Feuerbach.dist_circumcenter_excenter_sq},
copied from the author's Feuerbach development. Both
are also proved in Wang's Feuerbach repository and open Mathlib pull request
(Section~\ref{sec:related}).

\begin{theorem}[porism]\label{thm:porism}
Assume $0<r$, $r<R$ and $d^2=R(R-2r)$. For every $A\in\Gamma$ there are $B,C\in\Gamma$, with
$A,B,C$ pairwise distinct, such that the lines $AB$, $AC$ and $BC$ are tangent to~$\omega$.
\end{theorem}

This, and every Poncelet statement below, asserts tangency of the three side \emph{lines}. It does
not assert that $\omega$ is the incircle of $ABC$, that the points of tangency lie on the side
segments, or that the disk bounded by $\omega$ lies in the triangle. \emph{Informally}, that
classical reading follows under the hypotheses of Theorem~\ref{thm:porism}: they give
$d<R-r<R$ (\lean{Poncelet3.dist_center_lt}), so $I$ is inside $\Gamma$; by the classification of
points equidistant from the side lines and Euler's excentral formula used in the proof of
Theorem~\ref{thm:converse}, $I$ cannot be an excentre, so $\omega$ is the incircle. In the main files this argument
is not a formal output: the converse proof contains most of it, but the theorem it returns is only
Euler's relation. Version~2 proves it in the Layer-15 file, for $I$ inside $\Gamma$: three tangent side lines
make $\omega$ Mathlib's insphere of the triangle (\lean{Poncelet3.Classical.insphere_eq_of_isTangent},
Section~\ref{sec:l1415}).

\begin{lstlisting}
theorem exists_poncelet_triangle [Fact (Module.finrank ℝ V = 2)] (Γ ω : Sphere P)
    (hr : 0 < ω.radius) (hrR : ω.radius < Γ.radius)
    (hd : dist Γ.center ω.center ^ 2 = Γ.radius * (Γ.radius - 2 * ω.radius))
    {A : P} (hA : A ∈ Γ) :
    ∃ B C : P, B ∈ Γ ∧ C ∈ Γ ∧ A ≠ B ∧ A ≠ C ∧ B ≠ C ∧
      ω.IsTangent line[ℝ, A, B] ∧ ω.IsTangent line[ℝ, A, C] ∧ ω.IsTangent line[ℝ, B, C]
\end{lstlisting}

\paragraph{Why each hypothesis.}
Mathlib's spheres may have any real radius. Given $0<r$ and $d^2=R(R-2r)$, the hypothesis $r<R$ is
equivalent to $R>0$: if $R>0$ then $R-2r=d^2/R\ge0$, so $R>r$. It therefore only excludes
spheres of nonpositive radius, for which $R(R-2r)>0$ can hold too. The hypothesis $r>0$ excludes a
point-circle. We do \emph{not} assume that $\omega$ lies inside $\Gamma$, or that $A$ lies
outside $\omega$: both follow (\lean{Poncelet3.radius_lt_dist_center_of_mem} shows $r<|AI|$).

\paragraph{Proof.}
We use coordinates and polynomial certificates, not the classical synthetic argument. Take $I$ as
origin, $u=A-I$ and $v=Ju$ with $J$ the rotation by a right angle, and set $p=|u|^2$. The points
$B$, $C$ and the three points of tangency are given by explicit formulas in $u$, $v$, $r$, $R$,
$s=\sqrt{p-r^2}$ and the coordinates $X,Y$ of $O-I$. The facts ``on $\Gamma$'', ``on $\omega$'',
``perpendicular to the radius'', ``collinear'' and two non-degeneracy identities are 16
\code{linear\_combination} identities modulo five relations, with coefficients found by
Gr\"obner-basis reduction in a computer algebra system. Tangency then follows from a dimension
count: in the plane, a line through two distinct points of $\omega$'s tangent space at $T$ equals
that tangent space.

\begin{theorem}[converse]\label{thm:converse}
Suppose $A,B,C\in\Gamma$ are pairwise distinct, the lines $AB$, $AC$, $BC$ are tangent to
$\omega$, and $|IO|<R$. Then $d^2=R(R-2r)$. Hence, if $0<r<R$ and $|IO|<R$, every point of
$\Gamma$ is a vertex of such a triangle if and only if $d^2=R(R-2r)$.
\end{theorem}

\paragraph{Why $I$ must be inside $\Gamma$.}
The triangle $ABC$ spans the plane, so $\Gamma$ is its circumcircle, and $I$ is at distance $r$
from all three side lines. By Mathlib's
\lean{Affine.Simplex.exists_forall_dist_eq_iff_exists_excenterExists_and_eq_excenter}, $I$ is the
incentre or an excentre. In the incentre case, Euler's formula applies. In an excentre case,
$d^2=R(R+2r_i)\ge R^2$, contradicting $|IO|<R$. Without this hypothesis the first statement of
Theorem~\ref{thm:converse} is false. For example, for $A=(0,0)$, $B=(1,\tfrac1{10})$,
$C=(1,-\tfrac1{10})$ and $\omega$ the excircle opposite $A$, we have $r_A\approx0.110<R\approx0.505$,
all three side lines are tangent to $\omega$, and $d^2=R(R+2r_A)\ne R(R-2r_A)$.

\paragraph{Informal remark (not formalised).} In the ``every point'' form, and only there, the
inside-centre hypothesis can be dropped while keeping $0<r<R$. By the argument above, any one
triangle forces $d^2=R(R-2r)$ or the excentral alternative $d^2=R(R+2r)$. In the latter,
$d^2>R^2$, so $d>R$, and $(R+r)^2-d^2=r^2>0$, so $R<d<R+r$. The point of $\Gamma$ on the ray from
$O$ towards $I$ is at distance $d-R<r$ from $I$, hence strictly inside $\omega$, and no tangent
to $\omega$ passes through it. So the universal-vertex property rules out the excentral
alternative. This does \emph{not} remove the inside-centre hypothesis from the one-triangle
converse, nor from the all-or-nothing equivalence of Theorem~\ref{thm:construction}; the
excircle example above shows that it cannot be removed from the former.

\paragraph{The construction form.}
The classical statement is about a construction: draw the two tangents from $A$ to $\omega$, let
them meet $\Gamma$ again at $B$ and $C$; then $BC$ is tangent to $\omega$. Theorem~\ref{thm:porism}
only asserts that \emph{some} such $B,C$ exist. We therefore also prove the construction form.

\begin{definition}\label{def:closes}
\lean{Poncelet3.ConstructionCloses}~$\Gamma\ \omega\ A$ holds if (i) there are two distinct lines
$\ell_1,\ell_2\in\omega.\lean{tangentsFrom}\ A$ and points $B\in\Gamma\cap\ell_1$,
$C\in\Gamma\cap\ell_2$ different from $A$; and (ii) for \emph{all} such $\ell_1\ne\ell_2$, $B$, $C$,
the line $BC$ is tangent to~$\omega$.
\end{definition}

Part (i) prevents the predicate from holding vacuously, for instance when $A$ is inside $\omega$
or when a tangent from $A$ touches $\Gamma$ only at~$A$.

\begin{theorem}\label{thm:construction}
Assume $0<r<R$ and $|IO|<R$, and let $A\in\Gamma$. Then
\lean{ConstructionCloses}~$\Gamma\ \omega\ A$ holds if and only if $d^2=R(R-2r)$, and if and only if
\lean{ConstructionCloses}~$\Gamma\ \omega\ A'$ holds for every $A'\in\Gamma$. Moreover, if
$0<r<R$ and $d^2=R(R-2r)$, then for any two distinct tangents $\ell_1,\ell_2$ from $A$ and any
second intersections $B$, $C$, the line $BC$ is tangent (\lean{Poncelet3.poncelet_construction}).
\end{theorem}

The one new algebraic step is completeness of the tangency points: if $T\in\omega$ and $A$ lies
on the tangent at $T$, then $T$ is one of the two explicit tangency points. From it,
\lean{tangentsFrom} $A$ has exactly two elements, and each meets $\Gamma$ in exactly one point
other than~$A$.

\section{Do the statements mean what they say?}\label{sec:correct}

A reader of the statements above must trust several definitions. We list the lemmas that support
the intended reading.

\begin{lemma}[Carnot]\label{lem:carnot-correct}
For $p$ in the plane of $t$, $\lean{signedSideDist}\ t\ i\ p=\lean{signedInfDist}\ t\ i\ p$. The
vertex $A_i$ is at signed distance $h_i>0$ from side~$i$
(\lean{Carnot.signedSideDist_point_self}), and the incentre is at signed distance $r$ from every
side (\lean{Carnot.signedSideDist_incenter}).
\end{lemma}

Together with Theorem~\ref{thm:carnot-cos} ($d_i(O)=R\cos\alpha_i$) this pins down the sign
convention: the distance is positive on the side of the interior, and negative for the circumcentre
exactly when the opposite angle is obtuse. The circumcentre lies in the plane of $t$ by Mathlib's
\lean{Affine.Simplex.circumcenter_mem_affineSpan}, so Carnot's theorem could equally be stated with
Mathlib's \lean{signedInfDist}; the main files do not, but the Layer-14 file does
(\lean{Carnot.l14_carnot}, Section~\ref{sec:l1415}).

\begin{lemma}[Viviani]\label{lem:viv-unsigned}
For $p$ in the closed simplex, $0\le d_i(p)$ and
$d_i(p)=\operatorname{dist}(p,\ \pi_i(p))$, where $\pi_i$ is the orthogonal projection onto the
affine span of the $i$-th facet.
\end{lemma}

So the unsigned Viviani statement uses ordinary distances, computed from Mathlib's
\lean{Affine.Simplex.faceOpposite} and \lean{Affine.Simplex.orthogonalProjectionSpan}, and does not
depend on our choice of signed distance.

\begin{lemma}[De~Gua]\label{lem:degua-bridge}
In $\RR^3$, $\lean{triAreaSq}\ X\ Y\ Z=\lean{simplexVolSq}\,(X,Y,Z)$ (Lagrange's identity), and
Theorem~\ref{thm:degua3} is the case $m=2$ of Theorem~\ref{thm:degua}
(\lean{DeGua.de_gua_three_of_general}).
\end{lemma}

So the Gram-determinant volume agrees with the cross-product area in the one case where both are
familiar. A reader who is not convinced by Definition~\ref{def:volsq} can rely on
Theorem~\ref{thm:degua3}, whose definition of area is elementary.

\paragraph{Japanese theorem.}
Proposition~\ref{prop:bij} and Proposition~\ref{prop:cross} are correctness lemmas: they show that
the tree model, the combinatorial definition and the geometric definition of a triangulation all
describe the same objects, and that the triangles are the expected ones, $n-2$ in number
(\lean{Japanese.IsTriangulation.card_triangleSet}). \lean{Japanese.triInradius} returns $0$ for
collinear points, but under \code{hconv} no triangle used is degenerate.

\begin{lemma}[Poncelet]\label{lem:pon-correct}
Under the hypotheses of Theorem~\ref{thm:porism}, for $A\in\Gamma$:
\begin{enumerate}
\item[(a)] $A$ lies outside $\omega$ (\lean{Poncelet3.radius_lt_dist_center_of_mem});
\item[(b)] $\omega.\lean{tangentsFrom}\ A=\{AB,AC\}$ for the triangle of Theorem~\ref{thm:porism},
and $AB\ne AC$ (\lean{Poncelet3.tangentsFrom_eq_pair}, \lean{Poncelet3.ncard_tangentsFrom});
\item[(c)] for a tangent $\ell$ from $A$ with nonzero direction vector $v$, Mathlib's
$\Gamma.\lean{secondInter}\ A\ v$ is the unique point of $\Gamma\cap\ell$ other than $A$
(\lean{Poncelet3.secondInter_spec}).
\end{enumerate}
\end{lemma}

By (b), the $B$ and $C$ of Theorem~\ref{thm:porism} are exactly the points produced by the
classical construction, and by (c) that construction can be expressed with Mathlib's own
\lean{secondInter} (\lean{Poncelet3.poncelet_construction_secondInter}).

\section{Related work}\label{sec:related}

Our searches covered the pinned Mathlib source, the
current \texttt{docs/1000.yaml} of Mathlib's master branch (commit \texttt{6bd5e549},
30~September 2026), GitHub code and repository search, Mathlib issues and pull requests, and web
search. A search of the Lean Zulip chat on 30 September 2026 for ``Carnot'', ``Viviani'', ``De Gua'',
``Japanese theorem'' and ``Poncelet'' returned no relevant results.
The searches are not exhaustive, and the claims below are ``as far as we could determine''.
The per-project development notes record earlier rounds of these searches; where their early
conclusions (``no earlier Lean proof found'') predate the \texttt{lean-genius} findings below, this
section supersedes them. External code is cited at the commit we read (see the bibliography).

\paragraph{Mathlib.}
In both the pinned version and current master, \texttt{docs/1000.yaml} lists Viviani's theorem,
Carnot's theorem (inradius, circumradius), De~Gua's theorem, both Japanese theorems and Poncelet's
closure theorem with no declaration~\cite{Mathlib1000}. In the pinned \texttt{Mathlib/} source we found no
occurrence of these names and no Cauchy--Binet formula, and the module list of current master has
no Cauchy--Binet file. Mathlib does provide most of the vocabulary we use
(Section~\ref{sec:setting}), including \lean{signedInfDist}, whose \eqref{eq:bary} does most of the
work for Viviani and Carnot.

\paragraph{Euler's formula and Feuerbach's theorem.}
Wang's repository \texttt{wwylele/Feuerbach}~\cite{WangFeuerbach} proves Feuerbach's theorem in
Lean without \code{sorry}, and with it Euler's formula (\ext{dist_incenter_circumcenter}) and
Euler's inequality. His open Mathlib pull request \#44236~\cite{PR44236} adds, among others,
\ext{dist_incenter_circumcenter_sq} and \ext{dist_excenter_circumcenter_sq}. So Euler's
formula, on which our Poncelet development depends, is not new in Lean; our copy comes from an
independent development by the author. HOL~Light's \texttt{100/feuerbach.ml}~\cite{HOLLight} proves
Feuerbach's theorem in coordinates.

\paragraph{\texttt{lean-genius}.}
This public repository of Lean proofs~\cite{LeanGenius} contains several related files. We read
the five files cited below but did not build them; they contain no \code{sorry} and no declared
\code{axiom}.
\begin{itemize}
\item \texttt{DeGuaTheorem.lean} proves the three-dimensional De~Gua theorem for points of
\code{Fin 3 $\to$ $\mathbb{R}$}, with cross-product areas, by one \code{linear\_combination}. It is essentially our
Theorem~\ref{thm:degua3}, so the three-dimensional case is not new in Lean.
\item \texttt{VivianiTheorem.lean} proves Viviani's theorem for the unit equilateral triangle in
coordinates, with explicit line equations. \texttt{VivianiTheoremOQ01OQ02.lean} proves that for
any family of vectors whose sum is $(n+1)$ times a point $g$, the expression
$\sum_i(a-\langle R^{-1}(g-v_i),P\rangle)$ equals $(n+1)a$; its link to the facets of a regular
simplex is explained in comments but is not part of the statement. \texttt{VivianiTheoremOQ01OQ01.lean}
proves a converse for triangles in~$\mathbb C$.
\item \texttt{CarnotTheorem.lean} proves the trigonometric identity
$\cos A+\cos B+\cos C=1+4\sin\frac A2\sin\frac B2\sin\frac C2$ for reals with $A+B+C=\pi$. It does not
mention a triangle, circumcentre or inradius. Several follow-up files (\texttt{CarnotTheoremOQ01*})
develop related trigonometric identities and inequalities; we only read their headers.
\end{itemize}
Our Viviani theorem uses Mathlib's affine simplices and facet distances in every positive
dimension. Our De~Gua theorem uses Gram-determinant volumes for indexed vertex families in
arbitrary real inner product spaces, allowing degenerate configurations; it takes a point
\code{O : E} and a family \code{P : Fin (m+1) $\to$ E}, not an \lean{Affine.Simplex}.

\paragraph{Cauchy--Binet.}
The formula is not new as a formalisation, in Lean or elsewhere. In Lean: an open Mathlib pull request
(\#44209, labelled \texttt{awaiting-author} and \texttt{LLM-generated}) adds
\ext{Matrix.det_mul_eq_sum_det_submatrix}~\cite{PR44209}, and an earlier one (\#40473) was
closed without merging; \texttt{physlib} has a lemma \ext{cauchyBinet}~\cite{Physlib};
\texttt{rwst/Subspace-Theorems} has \ext{det_mul_eq_sum_det_submatrix}~\cite{Rwst}; and
\texttt{kim-em/hex-dev} has, for its own Mathlib-free matrix library, the ordered-column
expansion underlying Cauchy--Binet~\cite{HexDev}: \ext{det_columnSumMatrix_eq_sum_columnTuples}
expands the determinant as a sum over all ordered column tuples, repeated-column summands are shown
to vanish, and \ext{det_gramMatrix_eq_sum_columnTuples} specialises this to a Gram matrix. That
file's own header calls it the multilinearity form behind Cauchy--Binet; it does not state the
sum over column subsets of products of two minors. We have not checked the axioms of these. In
Isabelle/HOL, Divas\'on's AFP entry on the Smith normal form contains
\ext{Cauchy_Binet}~\cite{DivasonAFP}. In Coq, the CoqEAL library, built on MathComp, proves
\ext{BinetCauchy} in \path{theory/binetcauchy.v}~\cite{CoqEAL}; it is described by Cano, Cohen,
D\'en\`es, M\"ortberg and Siles~\cite{CanoEtAl}. Our proof is a self-contained alternative against
the pinned Mathlib, nothing more.

\paragraph{Other systems.}
We searched the Isabelle AFP (through its GitHub mirror), HOL~Light (a clone of the repository),
GeoCoq, and the web, for each of the five theorems.
\begin{itemize}
\item \emph{Isabelle AFP}: code search for ``Viviani'', ``Carnot'', ``Poncelet'', ``Japanese'',
``trirectangular'' and ``Chapple'' found nothing relevant. ``Gua'' is too common a substring to
search this way.
\item \emph{HOL~Light}: a case-insensitive search of the repository for ``viviani'', ``poncelet'',
``carnot'', ``japanese'', ``de gua'' and ``cauchy-binet'' found nothing.
\item \emph{GeoCoq}: code search for ``Viviani'', ``Carnot'', ``Poncelet'' and ``Japanese'' found
nothing, and the single hit for ``Gua'' is unrelated.
\item \emph{Mizar}: we only checked the article file names of the current MML for ``viviani'',
``carnot'', ``poncelet'', ``gua'' and ``binet'', and found none; we did not search article contents.
\item \emph{Web}: search for formalisations of each theorem returned only informal sources.
\end{itemize}
Apart from reading the CoqEAL file cited above, we did not search the Coq/Rocq ecosystem beyond GeoCoq. % TODO(Josh): check Rocq/Coq libraries beyond GeoCoq, and Mizar article contents

\section{Verification}\label{sec:verify}

All five projects were built with prebuilt Mathlib on a separate build machine; the machine that
prepared this paper did not run Lean. According to the development notes, Carnot, De~Gua and Poncelet were
rebuilt from an empty build directory (\code{rm -rf .lake/build; lake build}), and Viviani by a
``clean \code{lake build}''. The Japanese development was rebuilt from an empty build directory afterwards
(log \path{clean_build.log}: 8714 jobs, 25\,s, no warnings).
Build logs are stored for Viviani, De~Gua, Japanese and Poncelet; for Carnot only the notes
record the build. None of this was reproduced on the machine that prepared this paper.
Each build finished with no errors and no warnings. Wall-clock times: Carnot 15\,s, Viviani 6\,s,
De~Gua 10\,s, Japanese 25\,s (clean), Poncelet 30\,s.

\code{\#print axioms} was run on the headline theorems: 6 for Carnot, 8 for Viviani, 7 for
De~Gua, 29 for the Japanese development and 19 for Poncelet, 69 reports in all. The logs are kept
in each repository (\path{ax.log}, \path{scratch/axioms.out} for Viviani, or \path{axioms.txt}
for De~Gua). In each case the driver is a scratch file outside the library (\path{scratch/}).
Every report is \code{[propext, Classical.choice, Quot.sound]}. Every theorem in
Table~\ref{tab:thms} either has its own report or occurs in the statement or proof of a theorem
that has one. Where the development notes and the stored logs differ, the logs are authoritative. A grep over the library sources for
\code{sorry}, \code{admit}, \code{native\_decide}, \code{axiom}, \code{\#eval}, \code{run\_cmd},
\code{set\_option}, \code{macro}, \code{elab}, \code{syntax}, \code{notation} and
\code{import Lean} finds nothing. Each repository also has a \path{gate.sh} that rejects
forbidden tokens, builds the library, generates its own \code{\#print axioms} queries for every
declaration cited in this paper (plus those queried during development) and fails on any other
axiom. On the published commits it prints \code{PASS} with 12 (Carnot), 10 (Viviani), 10 (De~Gua),
51 (Japanese) and 22 (Poncelet) declarations (at the version-2 commit, 28, 28, 23, 81 and 53; Section~\ref{sec:l1415}). \path{tests/tamper.sh} plants six kinds of defect in
scratch copies (a \code{sorry}, a split \code{axiom}, a fake \code{\#eval} report, a macro, a
weakened statement, a wrapped extra axiom); in each repository the gate rejects all six, each with
a nonzero exit status and the expected failure message.

Each project also has a numerical sanity check in \texttt{code/}, written in Python with the same
definitions as the Lean code:
\begin{itemize}
\item \emph{Carnot}: 20\,000 random triangles, 14\,514 of them obtuse, checking both the sum and
$d_i=R\cos\alpha_i$; maximum relative error $4.5\cdot10^{-14}$.
\item \emph{Viviani}: 200 random simplices in each dimension $2$--$5$ with arbitrary affine
weights, and scaled regular simplices; maximum error $5.9\cdot10^{-12}$.
\item \emph{De~Gua}: the three-dimensional and general forms in dimensions up to 7, and
Cauchy--Binet including $k>n$ and $k=0$; maximum relative error $2\cdot10^{-12}$.
\item \emph{Japanese}: all triangulations of 200 random convex cyclic $n$-gons for $n=4,\dots,7$
(counts $2,5,14,42$). The tree model and the maximal non-crossing sets are equal as families for
$n\le8$, and crossing agrees with interleaving for 293\,088 chord pairs.
\item \emph{Poncelet}: 2\,000 random configurations; the explicit points lie on $\Gamma$ with all
three distances equal to $r$ (relative error $1.1\cdot10^{-12}$); with $d$ perturbed, the third
side misses $\omega$.
\end{itemize}

The step that matters most, comparing the formal statements with the informal theorems, cannot be
automated; Section~\ref{sec:correct} is meant to help with it.
The author has checked each statement in Table~\ref{tab:thms} against its informal counterpart by hand.

\subsection{Layers 14 and 15}\label{sec:l1415}

Version~2 adds two further checks. Their files are \path{<Lib>/L14.lean} and
\path{<Lib>/L15Classical.lean} in each repository (\path{Carnot}, \path{Viviani}, \path{DeGua},
\path{Japanese}, \path{Poncelet3}), at the new commit of each listed under Availability. Every theorem and lemma
in them is listed in that repository's \path{gate.sh}, which prints \code{PASS} with 28 (Carnot),
28 (Viviani), 23 (De~Gua), 81 (Japanese) and 53 (Poncelet) declarations, all with the three standard
axioms only, and \path{tests/tamper.sh} still rejects all six planted defects. The earlier library
files are byte-for-byte unchanged; only the root file of each library gains two import lines.
The published \path{gate.sh} used GNU \code{mktemp} options and failed on macOS (BSD tools) before
checking anything. The new commit uses a portable temporary file name, and the gate passes with BSD
tools. The \code{wrap} test of \path{tests/tamper.sh} still needs GNU \code{sed}.

\paragraph{Layer 14: statements in Mathlib's language, with witnesses.}
In \path{<Lib>/L14.lean}, 18 of the numbered results are restated with Mathlib-only types and
proved from the library theorems through definitional bridges: Theorems~\ref{thm:carnot}
and~\ref{thm:carnot-cos}, with Mathlib's \lean{signedInfDist} in place of
\lean{Carnot.signedSideDist} (\lean{Carnot.l14_carnot}); Theorems~\ref{thm:viv-weights}
and~\ref{thm:viv-heights}, Proposition~\ref{prop:equi} and Corollary~\ref{cor:viviani}, with
distances and heights written as \lean{Metric.infDist} to the facet span, and the classical
triangle form \lean{Viviani.l14_viviani_triangle}; Theorems~\ref{thm:degua3}, \ref{thm:degua}
and~\ref{thm:cb}, Lemma~\ref{lem:mdl} and Corollary~\ref{cor:proj}, with \lean{triAreaSq} and
\lean{simplexVolSq} unfolded to the cross product and to $\det$ of a Gram matrix divided by $(k!)^2$;
Lemma~\ref{lem:carnot-chord}, Theorem~\ref{thm:jq} and Theorem~\ref{thm:japanese}, the last with
\lean{IsGeomTriangulation} written out (\lean{Japanese.l14_japanese_polygon}); and
Theorems~\ref{thm:euler}, \ref{thm:porism}, \ref{thm:converse} and~\ref{thm:construction}, with
\lean{ConstructionCloses} unfolded into \lean{tangentsFrom} and \lean{IsTangent}. The elaborated type
of each restatement contains no constant from our libraries. No restatement turned out false or
needed an extra hypothesis. The Japanese restatements take an \lean{AffineIndependent} argument,
because a Mathlib triangle needs one; it follows from the orientation hypotheses and the conclusion
does not depend on it. For every restated result with hypotheses, a witness theorem proves all
hypotheses on an explicit configuration and proves that the conclusion is not degenerate: for
Carnot, the obtuse triangle $0,21,6+8i$ in $\mathbb C$, whose three signed distances $3R/5$,
$15R/17$, $-13R/85$ are distinct and of both signs (\lean{Carnot.l14w_carnot}); for Viviani, a point
with weights $(\frac12,\frac14,\frac14)$ in an equilateral triangle in $\RR^3$, whose distances are
positive and not all equal (\lean{Viviani.l14w_viviani}); for De~Gua, the corner $0,e_0,2e_1,3e_2$,
with pairwise distinct positive squared face areas $\frac{49}4,1,9,\frac94$, and for
Cauchy--Binet a $2\times3$ by $3\times2$ product with determinant $7\ne0$; for the Japanese
theorem, the quadrilateral $1,\frac{3+4i}5,-1,-i$ on the unit circle with its two distinct
triangulations; for Euler's formula, the triangle $0,4,3i$, with $O\ne I$; and for the porism, its
converse and the construction, circles with $d^2=\frac54=R(R-2r)$ and the explicit triangle
$0,4,3i$ with its three points of tangency (\lean{Poncelet3.l14w_poncelet}).
Lemma~\ref{lem:mdl} and Corollary~\ref{cor:proj} have no hypotheses and need no witness.
Not restated: Theorem~\ref{thm:jtree} and Propositions~\ref{prop:bij} and~\ref{prop:cross}, which
concern the tree model and the combinatorial definition (Theorem~\ref{thm:japanese} covers the
geometric statement); the correctness lemmas of Section~\ref{sec:correct}; and the auxiliary lemma
\lean{Carnot.carnot_poly}. These were left out to stay within the time box.

\paragraph{Layer 15: comparison with the classical statements.}
Another agent collected textbook statements of the six topics from the literature, with citations
(Johnson, Coxeter--Greitzer, Conant--Beyer, Fukagawa--Pedoe, Flatto and others). It cited them from
memory and did not check page numbers. An AI model given only these statements, with no access to our
code or to this paper, wrote them in Lean with Mathlib notions and definitions of its own: 22 theorems
and 3 well-formedness obligations, all with \code{sorry}. Its transcript and self-reported file list
were audited (a transcript audit in the verification repository; the writer's run passed). A comparer then copied each statement verbatim and either proved it from the libraries,
with no change to the library files, or classified it: 17 are \emph{covered} (proved), 6 are covered
only in a \emph{narrower} sense, and 2 are \emph{false as written} because the writer mis-stated
them. None was classified as out of scope alone, but one of the two false statements (the incentre
rectangle) is also outside the scope of this paper. The proofs are collected, with the writer's
definitions and statements reproduced verbatim, in \path{<Lib>/L15Classical.lean}, namespace
\code{<Lib>.Classical}, i.e.\ \code{Carnot.Classical} (and likewise \code{Viviani.Classical}, \code{DeGua.Classical},
\code{Japanese.Classical}, \code{Poncelet3.Classical}). All of them depend only on the three standard
axioms.

\emph{New proofs} (not in the main files):
\begin{itemize}
\item \lean{Carnot.Classical.signedDistLine_face_eq} (also in Viviani): for any positive-dimensional simplex and any point in its affine span, the writer's signed distance to a facet equals \lean{signedInfDist}.
\item \lean{Carnot.Classical.cov_TriSimplex1_b_acute}: for an acute triangle, Carnot's theorem with unsigned distances.
\item \lean{Carnot.Classical.cov_TriSimplex1_b_acute_inside}: the circumcentre of an acute triangle lies in its interior.
\item \lean{Viviani.Classical.height_two_eq}: the altitude of an equilateral triangle is $\frac{\sqrt3}2$ times the side; with it, Viviani with value $\frac{\sqrt3}2 s$ (open, closed, signed in the plane).
\item \lean{DeGua.Classical.cov_TriSimplex2_a_coord_part}: the coordinate form of three-dimensional De~Gua for \lean{simplexVolSq}.
\item \lean{Japanese.Classical.quad_signs}, \lean{Japanese.Classical.cov_TriSimplex4_b}: in the plane, strictly crossing diagonals together with positive orientation of $ABC$ give the other three orientation hypotheses of Theorem~\ref{thm:jq}, and the quadrilateral theorem holds for a convex quadrilateral (strictly crossing diagonals) with concyclic vertices, in any dimension.
\item \lean{Japanese.Classical.cov_TriSimplex4_a_obligation}: the triangles of a triangulation of a strictly convex polygon are non-degenerate.
\item \lean{Japanese.Classical.cov_TriSimplex4_b_obligation_part}: for pairwise distinct vertices, the four triangles of a convex quadrilateral are non-degenerate.
\item \lean{Poncelet3.Classical.cov_TriSimplex5_b}: the equality case of Euler's inequality, $R=2r$ if and only if the triangle is equilateral.
\item \lean{Poncelet3.Classical.insphere_eq_of_isTangent}: three tangent side lines and centre inside $\Gamma$ make $\omega$ the insphere of the triangle.
\item \lean{Poncelet3.Classical.cov_TriSimplex6_a}, \lean{Poncelet3.Classical.cov_TriSimplex6_a_construction}, \lean{Poncelet3.Classical.cov_TriSimplex6_b}, \lean{Poncelet3.Classical.cov_TriSimplex6_b_converse}, \lean{Poncelet3.Classical.cov_TriSimplex6_b_noclose}: the porism, the construction, the criterion and its two converses for triangles with vertices on $\Gamma$ and insphere $\omega$.
\end{itemize}

\emph{Narrower} (proved part, and what is missing):
\begin{itemize}
\item De~Gua in $\RR^3$, its coordinate form, the $n$-simplex form and its coordinate form (four statements): proved for \lean{simplexVolSq}; missing is the link to the writer's volume, Lebesgue measure of the hull in its span (and, for the $n$-simplex form, the indexing of facets by vertex sets).
\item The Japanese theorem for polygons: proved for maximal non-crossing chord sets and orientation convexity (\lean{Japanese.Classical.cov_TriSimplex4_a_part}); missing are the writer's dissections (cover with disjoint interiors) and side-line convexity.
\item Two tangents from every point of $\Gamma$: proved only under $d^2=R(R-2r)$ (\lean{Poncelet3.Classical.cov_TriSimplex6_a_obligation_part}); missing is the general case.
\end{itemize}

\emph{Writer statements false as written}:
\begin{itemize}
\item Non-degeneracy of the four triangles of a convex quadrilateral: the writer's convex quadrilateral allows coincident vertices (counterexample $A=B=0$, $C=D=1$ in $\RR$, \lean{Japanese.Classical.not_classical_TriSimplex4_b_obligation}); true for distinct vertices, as above.
\item The incentre rectangle: false for the same degenerate quadrilateral (\lean{Japanese.Classical.not_classical_TriSimplex4_b_rectangle}); the corrected statement is not proved and is outside this paper.
\end{itemize}
The narrower cases are the gaps already stated in Sections~\ref{sec:degua} and~\ref{sec:japanese} and
in Section~\ref{sec:limits}, and Lemma~\ref{lem:pon-correct}, which assumes $d^2=R(R-2r)$.

\section{Limitations and further work}\label{sec:limits}

\begin{itemize}
\item \emph{Duplication.} The Japanese development copies the Carnot files, and the Poncelet
development copies the Euler files. A shared library, or upstreaming to Mathlib, would remove this.
If PR~\#44236 is merged, the Euler copy should be replaced by Mathlib's lemmas.
\item \emph{Carnot with Mathlib's signed distance.} Theorem~\ref{thm:carnot} is stated directly with
\lean{signedInfDist} only in the Layer-14 file (\lean{Carnot.l14_carnot}), not in the main files.
\item \emph{Volumes.} \lean{simplexVolSq} is a definition, not linked to Mathlib's Euclidean
Hausdorff measure. For Viviani, we have not proved $h_i\operatorname{vol}(F_i)=n\operatorname{vol}(s)$,
which would turn ``equal heights'' into ``equal facet volumes''. The Layer-15 comparison
(Section~\ref{sec:l1415}) classifies the textbook De~Gua statements, written with Lebesgue measure
on the hull, as covered only in this narrower Gram-determinant sense.
\item \emph{Dissections.} The Japanese theorem is proved for maximal sets of chords with pairwise
disjoint open segments. We have not shown that every dissection of the polygon into triangles with
vertices among the $p_l$ arises this way; that needs an area or measure argument. For this reason the
Layer-15 comparison classifies the textbook polygon statement (triangulations as dissections, convexity
by side lines) as covered only in the narrower sense (Section~\ref{sec:l1415}).
\item \emph{Convex position.} \code{hconv} is phrased with orientations, not with Mathlib's
\lean{Convex} or \lean{convexHull}, and the polygon is indexed by $\mathbb N$ rather than
\code{Fin n}. For planar quadrilaterals, the Layer-15 lemma \lean{Japanese.Classical.quad_signs} derives the other three orientation hypotheses from strictly crossing diagonals and positive orientation of $ABC$; the full quadrilateral proof chooses this orientation in the non-degenerate case and treats degenerate cases separately.
\item \emph{Poncelet.} Only triangles are treated. Poncelet's closure theorem for $n$-gons, even for
$n=4$ (Fuss's relation), would need a different argument or new certificates. The theorems of
the main files assert side-line tangency only. The Layer-15 file (Section~\ref{sec:l1415}) proves,
for $0<r$ and $d<R-r$, the statements with an \lean{Affine.Triangle} whose vertices lie on $\Gamma$
and whose Mathlib insphere is $\omega$; that the points of tangency lie on the side segments is not
stated separately. That from every point of $\Gamma$ there are exactly two tangents to $\omega$, each
meeting $\Gamma$ again, is proved only under $d^2=R(R-2r)$ (Lemma~\ref{lem:pon-correct}).
\item \emph{Scope.} Lami's theorem, formalised separately by the author, is not part of this paper.
\end{itemize}

\section*{Acknowledgements and use of AI}
The formal proofs, the checking scripts and a draft of this paper were prepared with the help of
large-language-model AI assistants. The author compared each formal statement in
Table~\ref{tab:thms} with its informal counterpart by hand. The Layer-14 and Layer-15 checks of
Section~\ref{sec:l1415} (the classical statements, the Lean statements written from them, the
comparison proofs and the restatements) were produced by AI models. Every proof is checked by the Lean
kernel, and the author takes responsibility for all content.

\begin{thebibliography}{99}
\bibitem{CanoEtAl} G.~Cano, C.~Cohen, M.~D\'en\`es, A.~M\"ortberg and V.~Siles, Formalized linear
algebra over elementary divisor rings in Coq, \emph{Logical Methods in Computer Science} 12(2:7),
2016. \url{https://arxiv.org/abs/1601.07472}.
\bibitem{ConantBeyer} D.~R.~Conant and W.~A.~Beyer, Generalized Pythagorean theorem, \emph{The
American Mathematical Monthly} 81 (1974), no.~3, 262--265.
\bibitem{CoqEAL} CoqEAL contributors, \emph{CoqEAL: the Coq effective algebra library}, file
\path{theory/binetcauchy.v}, lemma \texttt{BinetCauchy}, \url{https://github.com/rocq-community/coqeal},
commit \texttt{8fd26cf4}, read 30 September 2026.
\bibitem{DivasonAFP} J.~Divas\'on, A verified algorithm for computing the Smith normal form of a
matrix, \emph{Archive of Formal Proofs}, May 2020, theory \texttt{Cauchy\_Binet},
\url{https://isa-afp.org/entries/Smith_Normal_Form.html}.
\bibitem{Flatto} L.~Flatto, \emph{Poncelet's Theorem}, American Mathematical Society, 2009.
\bibitem{FukagawaPedoe} H.~Fukagawa and D.~Pedoe, \emph{Japanese Temple Geometry Problems
(Sangaku)}, Charles Babbage Research Centre, Winnipeg, 1989.
\bibitem{HexDev} K.~Morrison, \emph{hex-dev}, file \path{HexDeterminant/CauchyBinet.lean}
(``Ordered column-tuple expansion underlying Cauchy-Binet''), \url{https://github.com/kim-em/hex-dev},
commit \texttt{f59485af} (file blob \texttt{4d2ab50a}), read 30 September 2026.
\bibitem{HOLLight} J.~Harrison, \texttt{100/feuerbach.ml}, HOL~Light repository,
\url{https://github.com/jrh13/hol-light}, commit \texttt{cba9198d}, read 30 September 2026.
\bibitem{LeanGenius} R.~J.~Walters, \emph{lean-genius}, files \path{proofs/Proofs/DeGuaTheorem.lean},
\path{VivianiTheorem.lean}, \path{VivianiTheoremOQ01OQ01.lean}, \path{VivianiTheoremOQ01OQ02.lean},
\path{CarnotTheorem.lean}, \url{https://github.com/rjwalters/lean-genius}, commit \texttt{8e708586}, read 30 September 2026.
\bibitem{Mathlib1000} The mathlib Community, \texttt{docs/1000.yaml}, Mathlib repository,
\url{https://github.com/leanprover-community/mathlib4}, pinned revision and master, accessed
30 September 2026.
\bibitem{PR44209} Mathlib pull request \#44209, ``feat(LinearAlgebra/Matrix): add the Cauchy-Binet
determinant formula'', \url{https://github.com/leanprover-community/mathlib4/pull/44209}, open as of
30 September 2026.
\bibitem{PR44236} W.~Wang, Mathlib pull request \#44236, ``feat(Geometry/Euclidean): prove
Feuerbach's theorem'', \url{https://github.com/leanprover-community/mathlib4/pull/44236}, open as of
30 September 2026.
\bibitem{Physlib} \emph{physlib}, file \path{QuantumInfo/ForMathlib/Majorization.lean},
\url{https://github.com/leanprover-community/physlib}, commit \texttt{d9108194}, read 30 September 2026.
\bibitem{Rwst} \emph{Subspace-Theorems}, file \path{ArithmeticHeights/CauchyBinet.lean},
\url{https://github.com/rwst/Subspace-Theorems}, commit \texttt{f10f915b}, read 30 September 2026.
\bibitem{WangFeuerbach} W.~Wang, \emph{Feuerbach}, \url{https://github.com/wwylele/Feuerbach},
commit \texttt{bdd66fa4}, read 30 September 2026.
\end{thebibliography}

\end{document}
